\documentclass[12pt,table,xcolor={dvipsnames}]{beamer}
\usetheme{Pittsburgh}
\usecolortheme{seagull}
%\usepackage[utf8]{inputenc}
\usepackage{fontspec}
\usepackage[brazilian]{babel}
\usepackage{amsmath}
\usepackage{listings}
\usepackage{multirow}
\usepackage{amsfonts}
\usepackage{amssymb}
\usepackage{graphicx}
\author{Pilhas e Filas Utilizando Vetores}
\title{Estruturas de Dados}
%\setbeamercovered{transparent}
\setbeamertemplate{navigation symbols}{}
%\logo{\includegraphics[scale=0.015]{Brasao_UFSC.png}\includegraphics[scale=0.2]{brasao_PPGCC.jpg}}
\institute{Departamento de Informática e de Estatística \\ Prof. Jean Everson Martina \\ Prof. Aldo von Wangenheim}
\date{2016.2}
\subject{}
\usebackgroundtemplate{\includegraphics[width=\paperwidth,
height=\paperheight]{../reusable_images/fundo_UFSC.png}}
\begin{document}

{
\usebackgroundtemplate{\includegraphics[width=\paperwidth,
height=\paperheight]{../reusable_images/fundo_capa.png}}
\begin{frame}
\titlepage
\includegraphics[scale=0.3]{../reusable_images/brasao_INE.png}
\end{frame}
}

\begin{frame}{Definição de Estruturas de Dados}
\begin{block}{Definição}
Estruturas de Dados é a disciplina que estuda as técnicas computacionais para a organização e manipulação eficiente de quaisquer quantidades de informação.
\end{block}
\end{frame}

\begin{frame}{Motivação de Estruturas de Dados}
\begin{itemize}
\item Ao estudar estruturas de dados teremos sempre este par:
\begin{itemize}
\item Um conjunto estruturado de informações usualmente uma classe de objetos ou um tipo de dados;
\item Um conjunto definido de operações sobre estes dados usualmente um conjunto de métodos ou funções.
\end{itemize}
\item Em um projeto de software sempre vamos nos preocupar de qual forma estão organizados os dados - qual a sua estrutura, e quais procedimentos que atuam sobre estes dados - que operações podem ser realizadas sobre eles.
\end{itemize}
\end{frame}

\begin{frame}{Conceito de Pilhas}
\begin{columns}
\column{.5\textwidth}
\begin{block}{Pilha}
É uma estrutura de dados cujo funcionamento é inspirado no conceito “natural” de empilhar, onde o primeiro a entrar é o último a sair (LIFO).
\end{block}
\column{.5\textwidth}
\begin{center}
\end{center}
\end{columns}
\end{frame}

\begin{frame}{Conceito de Pilhas}
\begin{columns}
\column{.5\textwidth}
\begin{block}{Pilha}
É uma estrutura de dados cujo funcionamento é inspirado no conceito “natural” de empilhar, onde o primeiro a entrar é o último a sair (LIFO).
\end{block}
\column{.5\textwidth}
\begin{center}
{\rowcolors{1}{SpringGreen}{GreenYellow}
\begin{tabular}{ |p{.5cm}| }
\hline
 \\ \hline
 \\ \hline
 \\ \hline
 \\ \hline
 \\ \hline
  \\ \hline
\end{tabular}
}
\end{center}
\end{columns}
\end{frame}

\begin{frame}{Conceito de Pilhas}
\begin{columns}
\column{.5\textwidth}
\begin{block}{Pilha}
É uma estrutura de dados cujo funcionamento é inspirado no conceito “natural” de empilhar, onde o primeiro a entrar é o último a sair (LIFO).
\end{block}
\column{.5\textwidth}
\begin{center}
{\rowcolors{1}{SpringGreen}{GreenYellow}
\begin{tabular}{ |p{.5cm}| }
\hline
 \\ \hline
 \\ \hline
 \\ \hline
 \\ \hline
 \\ \hline
20  \\ \hline
\end{tabular}
}
\end{center}
\end{columns}
\end{frame}

\begin{frame}{Conceito de Pilhas}
\begin{columns}
\column{.5\textwidth}
\begin{block}{Pilha}
É uma estrutura de dados cujo funcionamento é inspirado no conceito “natural” de empilhar, onde o primeiro a entrar é o último a sair (LIFO).
\end{block}
\column{.5\textwidth}
\begin{center}
{\rowcolors{1}{SpringGreen}{GreenYellow}
\begin{tabular}{ |p{.5cm}| }
\hline
 \\ \hline
 \\ \hline
 \\ \hline
 \\ \hline
55 \\ \hline
20  \\ \hline
\end{tabular}
}
\end{center}
\end{columns}
\end{frame}

\begin{frame}{Conceito de Pilhas}
\begin{columns}
\column{.5\textwidth}
\begin{block}{Pilha}
É uma estrutura de dados cujo funcionamento é inspirado no conceito “natural” de empilhar, onde o primeiro a entrar é o último a sair (LIFO).
\end{block}
\column{.5\textwidth}
\begin{center}
{\rowcolors{1}{SpringGreen}{GreenYellow}
\begin{tabular}{ |p{.5cm}| }
\hline
 \\ \hline
 \\ \hline
 \\ \hline
4 \\ \hline
55 \\ \hline
20  \\ \hline
\end{tabular}
}
\end{center}
\end{columns}
\end{frame}

\begin{frame}{Conceito de Pilhas}
\begin{columns}
\column{.5\textwidth}
\begin{block}{Pilha}
É uma estrutura de dados cujo funcionamento é inspirado no conceito “natural” de empilhar, onde o primeiro a entrar é o último a sair (LIFO).
\end{block}
\column{.5\textwidth}
\begin{center}
{\rowcolors{1}{SpringGreen}{GreenYellow}
\begin{tabular}{ |p{.5cm}| }
\hline
 \\ \hline
 \\ \hline
12 \\ \hline
4 \\ \hline
55 \\ \hline
20  \\ \hline
\end{tabular}
}
\end{center}
\end{columns}
\end{frame}

\begin{frame}{Conceito de Pilhas}
\begin{columns}
\column{.5\textwidth}
\begin{block}{Pilha}
É uma estrutura de dados cujo funcionamento é inspirado no conceito “natural” de empilhar, onde o primeiro a entrar é o último a sair (LIFO).
\end{block}
\column{.5\textwidth}
\begin{center}
{\rowcolors{1}{SpringGreen}{GreenYellow}
\begin{tabular}{ |p{.5cm}| }
\hline
 \\ \hline
89 \\ \hline
12 \\ \hline
4 \\ \hline
55 \\ \hline
20  \\ \hline
\end{tabular}
}
\end{center}
\end{columns}
\end{frame}

\begin{frame}{Conceito de Pilhas}
\begin{columns}
\column{.5\textwidth}
\begin{block}{Pilha}
É uma estrutura de dados cujo funcionamento é inspirado no conceito “natural” de empilhar, onde o primeiro a entrar é o último a sair (LIFO).
\end{block}
\column{.5\textwidth}
\begin{center}
{\rowcolors{1}{SpringGreen}{GreenYellow}
\begin{tabular}{ |p{.5cm}| }
\hline
24 \\ \hline
89 \\ \hline
12 \\ \hline
4 \\ \hline
55 \\ \hline
20  \\ \hline
\end{tabular}
}
\end{center}
\end{columns}
\end{frame}

\begin{frame}[fragile]{Pilhas Usando Vetores}
\begin{columns}
\column{.5\textwidth}
\begin{itemize}
\item Vetores possuem um espaço limitado para armazenar dados;
\item Necessitamos definir um espaço grande o suficiente para a nossa pilha;
\item Necessitamos de um indicador de qual elemento do vetor é o atual topo da pilha.
\end{itemize}
\column{.5\textwidth}
\begin{center}
\begin{tabular}{| p{.5cm} |p{.5cm}c }
\cline{1-1}
 \cellcolor{SpringGreen}  & &\\ \cline{1-1}
\cellcolor{GreenYellow}  & &\\ \cline{1-1}
\cellcolor{SpringGreen}  & &\\ \cline{1-1}
\cellcolor{GreenYellow}   & &\\ \cline{1-1}
\cellcolor{SpringGreen}  & &\\ \cline{1-2}
\cellcolor{GreenYellow}  & \multicolumn{1}{|p{.5cm}|} {int} & \\ \cline{1-2}
\end{tabular}
\end{center}
\end{columns}
\end{frame}

\begin{frame}[fragile]{Pilhas Usando Vetores}
\begin{columns}
\column{.5\textwidth}
\begin{itemize}
\item Vetores possuem um espaço limitado para armazenar dados;
\item Necessitamos definir um espaço grande o suficiente para a nossa pilha;
\item Necessitamos de um indicador de qual elemento do vetor é o atual topo da pilha.
\item {\color{red} Pilha Vazia!}
\end{itemize}
\column{.5\textwidth}
\begin{center}
\begin{tabular}{| p{.5cm} |p{.5cm}c }
\cline{1-1}
 \cellcolor{SpringGreen}  & &\\ \cline{1-1}
\cellcolor{GreenYellow}  & &\\ \cline{1-1}
\cellcolor{SpringGreen}  & &\\ \cline{1-1}
\cellcolor{GreenYellow}   & &\\ \cline{1-1}
\cellcolor{SpringGreen}  & &\\ \cline{1-2}
\cellcolor{GreenYellow}  & \multicolumn{1}{|p{.5cm}|} {-1} & \\ \cline{1-2}
\end{tabular}
\end{center}
\end{columns}
\end{frame}

\begin{frame}[fragile]{Pilhas Usando Vetores}
\begin{columns}
\column{.5\textwidth}
\begin{itemize}
\item Vetores possuem um espaço limitado para armazenar dados;
\item Necessitamos definir um espaço grande o suficiente para a nossa pilha;
\item Necessitamos de um indicador de qual elemento do vetor é o atual topo da pilha.
\end{itemize}
\column{.5\textwidth}
\begin{center}
\begin{tabular}{| p{.5cm} |p{.5cm}c }
\cline{1-1}
 \cellcolor{SpringGreen}  & &\\ \cline{1-1}
\cellcolor{GreenYellow}  & &\\ \cline{1-1}
\cellcolor{SpringGreen}  & &\\ \cline{1-1}
\cellcolor{GreenYellow}   & &\\ \cline{1-1}
\cellcolor{SpringGreen}  & &\\ \cline{1-2}
\cellcolor{GreenYellow} 20 & \multicolumn{1}{|p{.5cm}|} {0} & \\ \cline{1-2}
\end{tabular}
\end{center}
\end{columns}
\end{frame}

\begin{frame}[fragile]{Pilhas Usando Vetores}
\begin{columns}
\column{.5\textwidth}
\begin{itemize}
\item Vetores possuem um espaço limitado para armazenar dados;
\item Necessitamos definir um espaço grande o suficiente para a nossa pilha;
\item Necessitamos de um indicador de qual elemento do vetor é o atual topo da pilha.
\end{itemize}
\column{.5\textwidth}
\begin{center}
\begin{tabular}{| p{.5cm} |p{.5cm}c }
\cline{1-1}
 \cellcolor{SpringGreen}  & &\\ \cline{1-1}
\cellcolor{GreenYellow}  & &\\ \cline{1-1}
\cellcolor{SpringGreen}  & &\\ \cline{1-1}
\cellcolor{GreenYellow}   & &\\ \cline{1-1}
\cellcolor{SpringGreen} 55 & &\\ \cline{1-2}
\cellcolor{GreenYellow} 20 & \multicolumn{1}{|p{.5cm}|} {1} & \\ \cline{1-2}
\end{tabular}
\end{center}
\end{columns}
\end{frame}

\begin{frame}[fragile]{Pilhas Usando Vetores}
\begin{columns}
\column{.5\textwidth}
\begin{itemize}
\item Vetores possuem um espaço limitado para armazenar dados;
\item Necessitamos definir um espaço grande o suficiente para a nossa pilha;
\item Necessitamos de um indicador de qual elemento do vetor é o atual topo da pilha.
\end{itemize}
\column{.5\textwidth}
\begin{center}
\begin{tabular}{| p{.5cm} |p{.5cm}c }
\cline{1-1}
 \cellcolor{SpringGreen}  & &\\ \cline{1-1}
\cellcolor{GreenYellow}  & &\\ \cline{1-1}
\cellcolor{SpringGreen}  & &\\ \cline{1-1}
\cellcolor{GreenYellow}  4 & &\\ \cline{1-1}
\cellcolor{SpringGreen} 55 & &\\ \cline{1-2}
\cellcolor{GreenYellow} 20 & \multicolumn{1}{|p{.5cm}|} {2} & \\ \cline{1-2}
\end{tabular}
\end{center}
\end{columns}
\end{frame}

\begin{frame}[fragile]{Pilhas Usando Vetores}
\begin{columns}
\column{.5\textwidth}
\begin{itemize}
\item Vetores possuem um espaço limitado para armazenar dados;
\item Necessitamos definir um espaço grande o suficiente para a nossa pilha;
\item Necessitamos de um indicador de qual elemento do vetor é o atual topo da pilha.
\end{itemize}
\column{.5\textwidth}
\begin{center}
\begin{tabular}{| p{.5cm} |p{.5cm}c }
\cline{1-1}
 \cellcolor{SpringGreen}  & &\\ \cline{1-1}
\cellcolor{GreenYellow}  & &\\ \cline{1-1}
\cellcolor{SpringGreen}  12& &\\ \cline{1-1}
\cellcolor{GreenYellow}  4 & &\\ \cline{1-1}
\cellcolor{SpringGreen} 55 & &\\ \cline{1-2}
\cellcolor{GreenYellow} 20 & \multicolumn{1}{|p{.5cm}|} {3} & \\ \cline{1-2}
\end{tabular}
\end{center}
\end{columns}
\end{frame}

\begin{frame}[fragile]{Pilhas Usando Vetores}
\begin{columns}
\column{.5\textwidth}
\begin{itemize}
\item Vetores possuem um espaço limitado para armazenar dados;
\item Necessitamos definir um espaço grande o suficiente para a nossa pilha;
\item Necessitamos de um indicador de qual elemento do vetor é o atual topo da pilha.
\end{itemize}
\column{.5\textwidth}
\begin{center}
\begin{tabular}{| p{.5cm} |p{.5cm}c }
\cline{1-1}
 \cellcolor{SpringGreen}  & &\\ \cline{1-1}
\cellcolor{GreenYellow}  89& &\\ \cline{1-1}
\cellcolor{SpringGreen}  12& &\\ \cline{1-1}
\cellcolor{GreenYellow}  4 & &\\ \cline{1-1}
\cellcolor{SpringGreen} 55 & &\\ \cline{1-2}
\cellcolor{GreenYellow} 20 & \multicolumn{1}{|p{.5cm}|} {4} & \\ \cline{1-2}
\end{tabular}
\end{center}
\end{columns}
\end{frame}

\begin{frame}[fragile]{Pilhas Usando Vetores}
\begin{columns}
\column{.5\textwidth}
\begin{itemize}
\item Vetores possuem um espaço limitado para armazenar dados;
\item Necessitamos definir um espaço grande o suficiente para a nossa pilha;
\item Necessitamos de um indicador de qual elemento do vetor é o atual topo da pilha.
\end{itemize}
\column{.5\textwidth}
\begin{center}
\begin{tabular}{| p{.5cm} |p{.5cm}c }
\cline{1-1}
 \cellcolor{SpringGreen} 24 & &\\ \cline{1-1}
\cellcolor{GreenYellow}  89& &\\ \cline{1-1}
\cellcolor{SpringGreen}  12& &\\ \cline{1-1}
\cellcolor{GreenYellow}  4 & &\\ \cline{1-1}
\cellcolor{SpringGreen} 55 & &\\ \cline{1-2}
\cellcolor{GreenYellow} 20 & \multicolumn{1}{|p{.5cm}|} {5} & \\ \cline{1-2}
\end{tabular}
\end{center}
\end{columns}
\end{frame}

\begin{frame}[fragile]{Pilhas Usando Vetores}
\begin{columns}
\column{.5\textwidth}
\begin{itemize}
\item Vetores possuem um espaço limitado para armazenar dados;
\item Necessitamos definir um espaço grande o suficiente para a nossa pilha;
\item Necessitamos de um indicador de qual elemento do vetor é o atual topo da pilha.
\item {\color{red} Pilha Cheia!}
\end{itemize}
\column{.5\textwidth}
\begin{center}
\begin{tabular}{| p{.5cm} |p{.5cm}c }
\cline{1-1}
 \cellcolor{SpringGreen} 24 & &\\ \cline{1-1}
\cellcolor{GreenYellow}  89& &\\ \cline{1-1}
\cellcolor{SpringGreen}  12& &\\ \cline{1-1}
\cellcolor{GreenYellow}  4 & &\\ \cline{1-1}
\cellcolor{SpringGreen} 55 & &\\ \cline{1-2}
\cellcolor{GreenYellow} 20 & \multicolumn{1}{|p{.5cm}|} {5} & \\ \cline{1-2}
\end{tabular}
\end{center}
\end{columns}
\end{frame}

\begin{frame}[fragile]{Modelagem da Pilha}
\lstset{language=C++,
          keywordstyle=\color{blue}\ttfamily,
          stringstyle=\color{red}\ttfamily,
          commentstyle=\color{OliveGreen}\ttfamily,
          breaklines=true,
          basicstyle=\ttfamily\footnotesize
          }

\begin{itemize}
\item Aspecto Estrutural:
\begin{itemize}
\item Necessitamos de um vetor para armazenar as informações;
\item Necessitamos de um indicador da posição atual do topo da pilha;
\end{itemize}
\begin{lstlisting}
constantes MAXPILHA = 100;

classe Pilha {
 T _dados[MAXPILHA]; // Vetor Estatico
 inteiro _topo;        
};
\end{lstlisting}
\end{itemize}
\end{frame}

\begin{frame}[fragile]{Modelagem da Pilha}
\lstset{language=C++,
          keywordstyle=\color{blue}\ttfamily,
          stringstyle=\color{red}\ttfamily,
          commentstyle=\color{OliveGreen}\ttfamily,
          breaklines=true,
          basicstyle=\ttfamily\footnotesize
          }

\begin{itemize}
\item Aspecto Estrutural:
\begin{itemize}
\item Necessitamos de um vetor para armazenar as informações;
\item Necessitamos de um indicador da posição atual do topo da pilha;
\end{itemize}
\begin{lstlisting}
constantes MAXPILHA = 100;

classe Pilha {
 T* _dados; // _dados = new T[MAXPILHA];
 inteiro _topo;
 inteiro _tam = MAXPILHA;
};
\end{lstlisting}
\end{itemize}
\end{frame}

\begin{frame}[fragile]{Modelagem da Pilha}
\lstset{language=C++,
          keywordstyle=\color{blue}\ttfamily,
          stringstyle=\color{red}\ttfamily,
          commentstyle=\color{OliveGreen}\ttfamily,
          breaklines=true,
          basicstyle=\ttfamily\footnotesize
          }

\begin{itemize}
\item Aspecto Funcional:
\begin{itemize}
\item Temos que colocar e retirar dados da pilha;
\item Temos que testar se a pilha está vazia ou cheia;
\item Temos que inicializar a pilha.
\end{itemize}
\end{itemize}
\end{frame}

\begin{frame}[fragile]{Modelagem da Pilha}
\lstset{language=C++,
          keywordstyle=\color{blue}\ttfamily,
          stringstyle=\color{red}\ttfamily,
          commentstyle=\color{OliveGreen}\ttfamily,
          breaklines=true,
          basicstyle=\ttfamily\footnotesize
          }

\begin{itemize}
\item Inicializar ou limpar:
\begin{itemize}
\item Pilha();
\item Pilha(int tam);
\item limpaPilha();
\end{itemize}
\item Testar se a pilha está vazia ou cheia:
\begin{itemize}
\item bool pilhaCheia();
\item bool pilhaVazia();
\end{itemize}
\item Colocar e retirar dados da pilha:
\begin{itemize}
\item empilha(T dado);
\item T desempilha();
\item T topo();
\end{itemize}
\end{itemize}
\end{frame}

\begin{frame}[fragile]{Método \textit{Pilha()}}
\lstset{language=C++,
          keywordstyle=\color{blue}\ttfamily,
          stringstyle=\color{red}\ttfamily,
          commentstyle=\color{OliveGreen}\ttfamily,
          breaklines=true,
          basicstyle=\ttfamily\footnotesize
          }
\begin{lstlisting}
Pilha()
inicio
  _dados = new T[_tam];
  _topo <- -1;
fim;
\end{lstlisting}
\end{frame}

\begin{frame}[fragile]{Método \textit{Pilha(int tam)}}
\lstset{language=C++,
          keywordstyle=\color{blue}\ttfamily,
          stringstyle=\color{red}\ttfamily,
          commentstyle=\color{OliveGreen}\ttfamily,
          breaklines=true,
          basicstyle=\ttfamily\footnotesize
          }
\begin{lstlisting}
Pilha(int tam)
inicio
  _tam = tam;
  _dados = new T[tam];
  _topo <- -1;
fim;
\end{lstlisting}
\end{frame}

\begin{frame}[fragile]{Método \textit{limpaPilha()}}
\lstset{language=C++,
          keywordstyle=\color{blue}\ttfamily,
          stringstyle=\color{red}\ttfamily,
          commentstyle=\color{OliveGreen}\ttfamily,
          breaklines=true,
          basicstyle=\ttfamily\footnotesize
          }
\begin{lstlisting}
void limpaPilha()
inicio
  _topo <- -1;
fim;
\end{lstlisting}
\end{frame}

\begin{frame}[fragile]{Método \textit{pilhaCheia()}}
\lstset{language=C++,
          keywordstyle=\color{blue}\ttfamily,
          stringstyle=\color{red}\ttfamily,
          commentstyle=\color{OliveGreen}\ttfamily,
          breaklines=true,
          basicstyle=\ttfamily\footnotesize
          }
\begin{lstlisting}
bool pilhaCheia()
inicio
 SE (_topo = _tam - 1) ENTAO
  RETORNE(Verdadeiro);
 SENAO
  RETORNE(Falso);
fim
\end{lstlisting}
\end{frame}

\begin{frame}[fragile]{Método \textit{pilhaVazia()}}
\lstset{language=C++,
          keywordstyle=\color{blue}\ttfamily,
          stringstyle=\color{red}\ttfamily,
          commentstyle=\color{OliveGreen}\ttfamily,
          breaklines=true,
          basicstyle=\ttfamily\footnotesize
          }
\begin{lstlisting}
bool pilhaVazia()
inicio
SE (_topo = -1) ENTAO
 RETORNE(Verdadeiro)
SENAO
 RETORNE(Falso);
fim;
\end{lstlisting}
\end{frame}

\begin{frame}[fragile]{Método \textit{empilha(T dado)}}
\lstset{language=C++,
          keywordstyle=\color{blue}\ttfamily,
          stringstyle=\color{red}\ttfamily,
          commentstyle=\color{OliveGreen}\ttfamily,
          breaklines=true,
          basicstyle=\ttfamily\footnotesize
          }
\begin{lstlisting}
void empilha(T dado)
inicio
 SE (pilhaCheia) ENTAO
  JogueExcecao(ERROPILHACHEIA);
 SENAO
  _topo <- _topo + 1
  _dados[_topo] <- dado;
 FIM SE
fim;
\end{lstlisting}
\end{frame}

\begin{frame}[fragile]{Método \textit{T desempilha()}}
\lstset{language=C++,
          keywordstyle=\color{blue}\ttfamily,
          stringstyle=\color{red}\ttfamily,
          commentstyle=\color{OliveGreen}\ttfamily,
          breaklines=true,
          basicstyle=\ttfamily\footnotesize
          }
\begin{lstlisting}
T desempilha()
inicio
 SE (pilhaVazia) ENTAO
  JogueExcecao(ERROPILHAVAZIA);
 SENAO
  _topo <- _topo - 1;
  RETORNE(_dados[_topo+1]);
 FIM SE
fim;
\end{lstlisting}
\end{frame}

\begin{frame}[fragile]{Método \textit{T topo()}}
\lstset{language=C++,
          keywordstyle=\color{blue}\ttfamily,
          stringstyle=\color{red}\ttfamily,
          commentstyle=\color{OliveGreen}\ttfamily,
          breaklines=true,
          basicstyle=\ttfamily\footnotesize
          }
\begin{lstlisting}
T topo()
inicio
 SE (pilhaVazia) ENTAO
  JogueExcecao(ERROPILHAVAZIA);
 SENAO
  RETORNE(_dados[_topo]);
 FIM SE
fim;
\end{lstlisting}
\end{frame}

\begin{frame}[fragile]{Trabalho Pilha de Vetor}
\lstset{language=C++,
          keywordstyle=\color{blue}\ttfamily,
          stringstyle=\color{red}\ttfamily,
          commentstyle=\color{OliveGreen}\ttfamily,
          breaklines=true,
          basicstyle=\ttfamily\footnotesize
          }
\begin{itemize}
\item Implemente uma classe Pilha todas as operações vistas;
\item Implemente a pilha usando Templates;
\item Implemente a pilha com um numero de elementos variável definido na instanciação;
\item Use as melhores práticas de orientação a objetos;
\item Documente todas as classes, métodos e atributos;
\item Aplique os testes unitários disponíveis no moodle da disciplina para validar sua estrutura de dados;
\item Entregue até a data definida no moodle.
\end{itemize}
\end{frame}

\begin{frame}{Conceito de Filas}
\begin{columns}
\column{.5\textwidth}
\begin{block}{Fila}
É uma estrutura de dados cujo funcionamento é inspirado no conceito “natural” de enfileiramento, onde o primeiro a entrar é o primeiro a sair (FIFO).
\end{block}
\column{.5\textwidth}
\begin{center}
\end{center}
\end{columns}
\end{frame}

\begin{frame}{Conceito de Filas}
\begin{columns}
\column{.5\textwidth}
\begin{block}{Fila}
É uma estrutura de dados cujo funcionamento é inspirado no conceito “natural” de enfileiramento, onde o primeiro a entrar é o primeiro a sair (FIFO).
\end{block}
\column{.5\textwidth}
\begin{center}
{\rowcolors{1}{SpringGreen}{GreenYellow}
\begin{tabular}{ |p{.5cm}|p{.5cm}|p{.5cm}|p{.5cm}|p{.5cm}|p{.5cm}| }
\hline
 & & & & &\\ \hline
\end{tabular}
}
\end{center}
\end{columns}
\end{frame}

\begin{frame}{Conceito de Filas}
\begin{columns}
\column{.5\textwidth}
\begin{block}{Fila}
É uma estrutura de dados cujo funcionamento é inspirado no conceito “natural” de enfileiramento, onde o primeiro a entrar é o primeiro a sair (FIFO).
\end{block}
\column{.5\textwidth}
\begin{center}
{\rowcolors{1}{SpringGreen}{GreenYellow}
\begin{tabular}{ |p{.5cm}|p{.5cm}|p{.5cm}|p{.5cm}|p{.5cm}|p{.5cm}| }
\hline
 &  &  &  &  & 20 \\ \hline
\end{tabular}
}
\end{center}
\end{columns}
\end{frame}

\begin{frame}{Conceito de Filas}
\begin{columns}
\column{.5\textwidth}
\begin{block}{Fila}
É uma estrutura de dados cujo funcionamento é inspirado no conceito “natural” de enfileiramento, onde o primeiro a entrar é o primeiro a sair (FIFO).
\end{block}
\column{.5\textwidth}
\begin{center}
{\rowcolors{1}{SpringGreen}{GreenYellow}
\begin{tabular}{ |p{.5cm}|p{.5cm}|p{.5cm}|p{.5cm}|p{.5cm}|p{.5cm}| }
\hline
 &  &  &  & 55 & 20 \\ \hline
\end{tabular}
}
\end{center}
\end{columns}
\end{frame}

\begin{frame}{Conceito de Filas}
\begin{columns}
\column{.5\textwidth}
\begin{block}{Fila}
É uma estrutura de dados cujo funcionamento é inspirado no conceito “natural” de enfileiramento, onde o primeiro a entrar é o primeiro a sair (FIFO).
\end{block}
\column{.5\textwidth}
\begin{center}
{\rowcolors{1}{SpringGreen}{GreenYellow}
\begin{tabular}{ |p{.5cm}|p{.5cm}|p{.5cm}|p{.5cm}|p{.5cm}|p{.5cm}| }
\hline
 &  &  & 4 & 55 & 20 \\ \hline
\end{tabular}
}
\end{center}
\end{columns}
\end{frame}

\begin{frame}{Conceito de Filas}
\begin{columns}
\column{.5\textwidth}
\begin{block}{Fila}
É uma estrutura de dados cujo funcionamento é inspirado no conceito “natural” de enfileiramento, onde o primeiro a entrar é o primeiro a sair (FIFO).
\end{block}
\column{.5\textwidth}
\begin{center}
{\rowcolors{1}{SpringGreen}{GreenYellow}
\begin{tabular}{ |p{.5cm}|p{.5cm}|p{.5cm}|p{.5cm}|p{.5cm}|p{.5cm}| }
\hline
 &  & 12 & 4 & 55 & 20 \\ \hline
\end{tabular}
}
\end{center}
\end{columns}
\end{frame}

\begin{frame}{Conceito de Filas}
\begin{columns}
\column{.5\textwidth}
\begin{block}{Fila}
É uma estrutura de dados cujo funcionamento é inspirado no conceito “natural” de enfileiramento, onde o primeiro a entrar é o primeiro a sair (FIFO).
\end{block}
\column{.5\textwidth}
\begin{center}
{\rowcolors{1}{SpringGreen}{GreenYellow}
\begin{tabular}{ |p{.5cm}|p{.5cm}|p{.5cm}|p{.5cm}|p{.5cm}|p{.5cm}| }
\hline
 & 89 & 12 & 4 & 55 & 20 \\ \hline
\end{tabular}
}
\end{center}
\end{columns}
\end{frame}

\begin{frame}{Conceito de Filas}
\begin{columns}
\column{.5\textwidth}
\begin{block}{Fila}
É uma estrutura de dados cujo funcionamento é inspirado no conceito “natural” de enfileiramento, onde o primeiro a entrar é o primeiro a sair (FIFO).
\end{block}
\column{.5\textwidth}
\begin{center}
{\rowcolors{1}{SpringGreen}{GreenYellow}
\begin{tabular}{ |p{.5cm}|p{.5cm}|p{.5cm}|p{.5cm}|p{.5cm}|p{.5cm}| }
\hline
24 & 89 & 12 & 4 & 55 & 20 \\ \hline
\end{tabular}
}
\end{center}
\end{columns}
\end{frame}

\begin{frame}{Conceito de Filas}
\begin{columns}
\column{.5\textwidth}
\begin{block}{Fila}
É uma estrutura de dados cujo funcionamento é inspirado no conceito “natural” de enfileiramento, onde o primeiro a entrar é o primeiro a sair (FIFO).
\end{block}
\column{.5\textwidth}
\begin{center}
{\rowcolors{1}{SpringGreen}{GreenYellow}
\begin{tabular}{ |p{.5cm}|p{.5cm}|p{.5cm}|p{.5cm}|p{.5cm}|p{.5cm}| }
\hline
& 24 & 89 & 12 & 4 & 55 \\ \hline
\end{tabular}
}
\end{center}
\end{columns}
\end{frame}

\begin{frame}{Conceito de Filas}
\begin{columns}
\column{.5\textwidth}
\begin{block}{Fila}
É uma estrutura de dados cujo funcionamento é inspirado no conceito “natural” de enfileiramento, onde o primeiro a entrar é o primeiro a sair (FIFO).
\end{block}
\column{.5\textwidth}
\begin{center}
{\rowcolors{1}{SpringGreen}{GreenYellow}
\begin{tabular}{ |p{.5cm}|p{.5cm}|p{.5cm}|p{.5cm}|p{.5cm}|p{.5cm}| }
\hline
& & 24 & 89 & 12 & 4  \\ \hline
\end{tabular}
}
\end{center}
\end{columns}
\end{frame}

\begin{frame}{Conceito de Filas}
\begin{columns}
\column{.5\textwidth}
\begin{block}{Fila}
É uma estrutura de dados cujo funcionamento é inspirado no conceito “natural” de enfileiramento, onde o primeiro a entrar é o primeiro a sair (FIFO).
\end{block}
\column{.5\textwidth}
\begin{center}
{\rowcolors{1}{SpringGreen}{GreenYellow}
\begin{tabular}{ |p{.5cm}|p{.5cm}|p{.5cm}|p{.5cm}|p{.5cm}|p{.5cm}| }
\hline
& & 24 & 89 & 12 & 4  \\ \hline
\end{tabular}
}
\end{center}
\begin{itemize}
\item Duas operações:
\begin{itemize}
\item {\color{red}Incluir no Fim};
\item {\color{red}Retirar do Inicio};
\end{itemize}
\end{itemize}
\end{columns}
\end{frame}

\begin{frame}{Conceito de Filas}
\begin{itemize}
\item É importante para gerência de dados/processos por ordem cronológica:
\begin{itemize}
\item Fila de impressão em uma impressora de rede;
\item Fila de pedidos de uma expedição ou tele-entrega;
\end{itemize}
\item É importante para simulação de processos sequenciais:
\begin{itemize}
\item chão de fábrica: fila de camisetas a serem estampadas;
\item comércio: simulação de fluxo de um caixa de supermercado;
\item tráfego: simulação de um cruzamento com um semáforo.
\end{itemize}
\end{itemize}
\end{frame}

\begin{frame}[fragile]{Filas Usando Vetores}
\begin{columns}
\column{.5\textwidth}
\begin{itemize}
\item Vetores possuem um espaço limitado para armazenar dados;
\item Necessitamos definir um espaço grande o suficiente para a nossa pilha;
\item Necessitamos de um indicador de qual elemento do vetor é o atual topo da pilha.
\item Incluímos sempre no fim.
\end{itemize}
\column{.5\textwidth}
\begin{center}
\begin{tabular}{| p{.5cm} |p{.5cm}c }
\cline{1-1}
 \cellcolor{SpringGreen}  & &\\ \cline{1-1}
\cellcolor{GreenYellow}  & &\\ \cline{1-1}
\cellcolor{SpringGreen}  & &\\ \cline{1-1}
\cellcolor{GreenYellow}   & &\\ \cline{1-1}
\cellcolor{SpringGreen}  & &\\ \cline{1-2}
\cellcolor{GreenYellow}  & \multicolumn{1}{|p{.5cm}|} {int} & \\ \cline{1-2}
\end{tabular}
\end{center}
\end{columns}
\end{frame}

\begin{frame}[fragile]{Filas Usando Vetores}
\begin{columns}
\column{.5\textwidth}
\begin{itemize}
\item Vetores possuem um espaço limitado para armazenar dados;
\item Necessitamos definir um espaço grande o suficiente para a nossa pilha;
\item Necessitamos de um indicador de qual elemento do vetor é o atual topo da pilha.
\item Incluímos sempre no fim.
\item {\color{red}Fila Vazia!}
\end{itemize}
\column{.5\textwidth}
\begin{center}
\begin{tabular}{| p{.5cm} |p{.5cm}c }
\cline{1-1}
 \cellcolor{SpringGreen}  & &\\ \cline{1-1}
\cellcolor{GreenYellow}  & &\\ \cline{1-1}
\cellcolor{SpringGreen}  & &\\ \cline{1-1}
\cellcolor{GreenYellow}   & &\\ \cline{1-1}
\cellcolor{SpringGreen}  & &\\ \cline{1-2}
\cellcolor{GreenYellow}  & \multicolumn{1}{|p{.5cm}|} {-1} & \\ \cline{1-2}
\end{tabular}
\end{center}
\end{columns}
\end{frame}

\begin{frame}[fragile]{Filas Usando Vetores}
\begin{columns}
\column{.5\textwidth}
\begin{itemize}
\item Vetores possuem um espaço limitado para armazenar dados;
\item Necessitamos definir um espaço grande o suficiente para a nossa pilha;
\item Necessitamos de um indicador de qual elemento do vetor é o atual topo da pilha.
\item Incluímos sempre no fim.
\end{itemize}
\column{.5\textwidth}
\begin{center}
\begin{tabular}{| p{.5cm} |p{.5cm}c }
\cline{1-1}
 \cellcolor{SpringGreen}  & &\\ \cline{1-1}
\cellcolor{GreenYellow}  & &\\ \cline{1-1}
\cellcolor{SpringGreen}  & &\\ \cline{1-1}
\cellcolor{GreenYellow}   & &\\ \cline{1-1}
\cellcolor{SpringGreen}  & &\\ \cline{1-2}
\cellcolor{GreenYellow} 20 & \multicolumn{1}{|p{.5cm}|} {0} & \\ \cline{1-2}
\end{tabular}
\end{center}
\end{columns}
\end{frame}

\begin{frame}[fragile]{Filas Usando Vetores}
\begin{columns}
\column{.5\textwidth}
\begin{itemize}
\item Vetores possuem um espaço limitado para armazenar dados;
\item Necessitamos definir um espaço grande o suficiente para a nossa pilha;
\item Necessitamos de um indicador de qual elemento do vetor é o atual topo da pilha.
\item Incluímos sempre no fim.
\end{itemize}
\column{.5\textwidth}
\begin{center}
\begin{tabular}{| p{.5cm} |p{.5cm}c }
\cline{1-1}
 \cellcolor{SpringGreen}  & &\\ \cline{1-1}
\cellcolor{GreenYellow}  & &\\ \cline{1-1}
\cellcolor{SpringGreen}  & &\\ \cline{1-1}
\cellcolor{GreenYellow}   & &\\ \cline{1-1}
\cellcolor{SpringGreen} 55 & &\\ \cline{1-2}
\cellcolor{GreenYellow} 20 & \multicolumn{1}{|p{.5cm}|} {1} & \\ \cline{1-2}
\end{tabular}
\end{center}
\end{columns}
\end{frame}

\begin{frame}[fragile]{Filas Usando Vetores}
\begin{columns}
\column{.5\textwidth}
\begin{itemize}
\item Vetores possuem um espaço limitado para armazenar dados;
\item Necessitamos definir um espaço grande o suficiente para a nossa pilha;
\item Necessitamos de um indicador de qual elemento do vetor é o atual topo da pilha.
\item Incluímos sempre no fim.
\end{itemize}
\column{.5\textwidth}
\begin{center}
\begin{tabular}{| p{.5cm} |p{.5cm}c }
\cline{1-1}
 \cellcolor{SpringGreen}  & &\\ \cline{1-1}
\cellcolor{GreenYellow}  & &\\ \cline{1-1}
\cellcolor{SpringGreen}  & &\\ \cline{1-1}
\cellcolor{GreenYellow}  4 & &\\ \cline{1-1}
\cellcolor{SpringGreen} 55 & &\\ \cline{1-2}
\cellcolor{GreenYellow} 20 & \multicolumn{1}{|p{.5cm}|} {2} & \\ \cline{1-2}
\end{tabular}
\end{center}
\end{columns}
\end{frame}

\begin{frame}[fragile]{Filas Usando Vetores}
\begin{columns}
\column{.5\textwidth}
\begin{itemize}
\item Vetores possuem um espaço limitado para armazenar dados;
\item Necessitamos definir um espaço grande o suficiente para a nossa pilha;
\item Necessitamos de um indicador de qual elemento do vetor é o atual topo da pilha.
\item Incluímos sempre no fim.
\end{itemize}
\column{.5\textwidth}
\begin{center}
\begin{tabular}{| p{.5cm} |p{.5cm}c }
\cline{1-1}
 \cellcolor{SpringGreen}  & &\\ \cline{1-1}
\cellcolor{GreenYellow}  & &\\ \cline{1-1}
\cellcolor{SpringGreen}  12& &\\ \cline{1-1}
\cellcolor{GreenYellow}  4 & &\\ \cline{1-1}
\cellcolor{SpringGreen} 55 & &\\ \cline{1-2}
\cellcolor{GreenYellow} 20 & \multicolumn{1}{|p{.5cm}|} {3} & \\ \cline{1-2}
\end{tabular}
\end{center}
\end{columns}
\end{frame}

\begin{frame}[fragile]{Filas Usando Vetores}
\begin{columns}
\column{.5\textwidth}
\begin{itemize}
\item Vetores possuem um espaço limitado para armazenar dados;
\item Necessitamos definir um espaço grande o suficiente para a nossa pilha;
\item Necessitamos de um indicador de qual elemento do vetor é o atual topo da pilha.
\item Incluímos sempre no fim.
\end{itemize}
\column{.5\textwidth}
\begin{center}
\begin{tabular}{| p{.5cm} |p{.5cm}c }
\cline{1-1}
 \cellcolor{SpringGreen}  & &\\ \cline{1-1}
\cellcolor{GreenYellow}  89& &\\ \cline{1-1}
\cellcolor{SpringGreen}  12& &\\ \cline{1-1}
\cellcolor{GreenYellow}  4 & &\\ \cline{1-1}
\cellcolor{SpringGreen} 55 & &\\ \cline{1-2}
\cellcolor{GreenYellow} 20 & \multicolumn{1}{|p{.5cm}|} {4} & \\ \cline{1-2}
\end{tabular}
\end{center}
\end{columns}
\end{frame}

\begin{frame}[fragile]{Filas Usando Vetores}
\begin{columns}
\column{.5\textwidth}
\begin{itemize}
\item Vetores possuem um espaço limitado para armazenar dados;
\item Necessitamos definir um espaço grande o suficiente para a nossa pilha;
\item Necessitamos de um indicador de qual elemento do vetor é o atual topo da pilha.
\item Incluímos sempre no fim.
\end{itemize}
\column{.5\textwidth}
\begin{center}
\begin{tabular}{| p{.5cm} |p{.5cm}c }
\cline{1-1}
 \cellcolor{SpringGreen} 24 & &\\ \cline{1-1}
\cellcolor{GreenYellow}  89& &\\ \cline{1-1}
\cellcolor{SpringGreen}  12& &\\ \cline{1-1}
\cellcolor{GreenYellow}  4 & &\\ \cline{1-1}
\cellcolor{SpringGreen} 55 & &\\ \cline{1-2}
\cellcolor{GreenYellow} 20 & \multicolumn{1}{|p{.5cm}|} {5} & \\ \cline{1-2}
\end{tabular}
\end{center}
\end{columns}
\end{frame}

\begin{frame}[fragile]{Filas Usando Vetores}
\begin{columns}
\column{.5\textwidth}
\begin{itemize}
\item Vetores possuem um espaço limitado para armazenar dados;
\item Necessitamos definir um espaço grande o suficiente para a nossa pilha;
\item Necessitamos de um indicador de qual elemento do vetor é o atual topo da pilha.
\item Incluímos sempre no fim.
\item {\color{red}Fila Cheia!}
\end{itemize}
\column{.5\textwidth}
\begin{center}
\begin{tabular}{| p{.5cm} |p{.5cm}c }
\cline{1-1}
 \cellcolor{SpringGreen} 24 & &\\ \cline{1-1}
\cellcolor{GreenYellow}  89& &\\ \cline{1-1}
\cellcolor{SpringGreen}  12& &\\ \cline{1-1}
\cellcolor{GreenYellow}  4 & &\\ \cline{1-1}
\cellcolor{SpringGreen} 55 & &\\ \cline{1-2}
\cellcolor{GreenYellow} 20 & \multicolumn{1}{|p{.5cm}|} {5} & \\ \cline{1-2}
\end{tabular}
\end{center}
\end{columns}
\end{frame}

\begin{frame}[fragile]{Modelagem da Fila}
\lstset{language=C++,
          keywordstyle=\color{blue}\ttfamily,
          stringstyle=\color{red}\ttfamily,
          commentstyle=\color{OliveGreen}\ttfamily,
          breaklines=true,
          basicstyle=\ttfamily\footnotesize
          }

\begin{itemize}
\item Aspecto Estrutural:
\begin{itemize}
\item Necessitamos de um vetor para armazenar as informações;
\item Necessitamos de um indicador da posição atual do fim da fila;
\end{itemize}
\begin{lstlisting}
constantes MAXFILA = 100;

classe Fila {
 T* _dados; // _dados = new T[MAXFILA];
 inteiro _fim;
 inteiro _tam = MAXFILA;
};
\end{lstlisting}
\end{itemize}
\end{frame}

\begin{frame}[fragile]{Modelagem da Fila}
\lstset{language=C++,
          keywordstyle=\color{blue}\ttfamily,
          stringstyle=\color{red}\ttfamily,
          commentstyle=\color{OliveGreen}\ttfamily,
          breaklines=true,
          basicstyle=\ttfamily\footnotesize
          }

\begin{itemize}
\item Aspecto Funcional:
\begin{itemize}
\item Temos que colocar e retirar dados da fila;
\item Temos que testar se a fila está vazia ou cheia;
\item Temos que inicializar a fila.
\end{itemize}
\end{itemize}
\end{frame}

\begin{frame}[fragile]{Modelagem da Fila}
\lstset{language=C++,
          keywordstyle=\color{blue}\ttfamily,
          stringstyle=\color{red}\ttfamily,
          commentstyle=\color{OliveGreen}\ttfamily,
          breaklines=true,
          basicstyle=\ttfamily\footnotesize
          }

\begin{itemize}
\item Inicializar ou limpar:
\begin{itemize}
\item Fila();
\item Fila(int tam);
\item limpaFila();
\end{itemize}
\item Testar se a fila está vazia ou cheia:
\begin{itemize}
\item bool filaCheia();
\item bool filaVazia();
\end{itemize}
\item Colocar e retirar dados da pilha:
\begin{itemize}
\item insere(T dado);
\item T retira();
\item T ultimo();
\end{itemize}
\end{itemize}
\end{frame}

\begin{frame}[fragile]{Método \textit{T retira()}}
\begin{columns}
\column{.5\textwidth}
\begin{itemize}
\item Testamos se há elementos;
\item Decrementamos o fim da fila (último);
\item Salvamos o primeiro elemento em variável auxiliar;
\item Empurramos tudo para a frente.
\end{itemize}
\column{.5\textwidth}
\begin{center}
\begin{tabular}{| p{.5cm} |p{.5cm}c }
\cline{1-1}
 \cellcolor{SpringGreen} 24 & &\\ \cline{1-1}
\cellcolor{GreenYellow}  89& &\\ \cline{1-1}
\cellcolor{SpringGreen}  12& &\\ \cline{1-1}
\cellcolor{GreenYellow}  4 & &\\ \cline{1-1}
\cellcolor{SpringGreen} 55 & &\\ \cline{1-2}
\cellcolor{GreenYellow} 20 & \multicolumn{1}{|p{.5cm}|} {5} & \\ \cline{1-2}
\end{tabular}
\end{center}
\end{columns}
\end{frame}

\begin{frame}[fragile]{Método \textit{T retira()}}
\begin{columns}
\column{.5\textwidth}
\begin{itemize}
\item Testamos se há elementos;
\item Decrementamos o fim da fila (último);
\item Salvamos o primeiro elemento em variável auxiliar;
\item Empurramos tudo para a frente.
\end{itemize}
\column{.5\textwidth}
\begin{center}
\begin{tabular}{| p{.5cm} |p{.5cm}c }
\cline{1-1}
 \cellcolor{SpringGreen} 24 & &\\ \cline{1-1}
\cellcolor{GreenYellow}  89& &\\ \cline{1-1}
\cellcolor{SpringGreen}  12& &\\ \cline{1-1}
\cellcolor{GreenYellow}  4 & &\\ \cline{1-1}
\cellcolor{SpringGreen} 55 & &\\ \cline{1-2}
\cellcolor{GreenYellow} 20 & \multicolumn{1}{|p{.5cm}|} {5} & \\ \cline{1-2}
\end{tabular}
\begin{tabular}{| p{.5cm} | }
\hline
\cellcolor{Mahogany} {T} \\\hline
\end{tabular}
\end{center}
\end{columns}
\end{frame}

\begin{frame}[fragile]{Método \textit{T retira()}}
\begin{columns}
\column{.5\textwidth}
\begin{itemize}
\item Testamos se há elementos;
\item Decrementamos o fim da fila (último);
\item Salvamos o primeiro elemento em variável auxiliar;
\item Empurramos tudo para a frente.
\end{itemize}
\column{.5\textwidth}
\begin{center}
\begin{tabular}{| p{.5cm} |p{.5cm}c }
\cline{1-1}
 \cellcolor{SpringGreen} 24 & &\\ \cline{1-1}
\cellcolor{GreenYellow}  89& &\\ \cline{1-1}
\cellcolor{SpringGreen}  12& &\\ \cline{1-1}
\cellcolor{GreenYellow}  4 & &\\ \cline{1-1}
\cellcolor{SpringGreen} 55 & &\\ \cline{1-2}
\cellcolor{GreenYellow}  & \multicolumn{1}{|p{.5cm}|} {4} & \\ \cline{1-2}
\end{tabular}
\begin{tabular}{| p{.5cm} | }
\hline
\cellcolor{Mahogany} {T} \\\hline
\end{tabular}
\end{center}
\end{columns}
\end{frame}

\begin{frame}[fragile]{Método \textit{T retira()}}
\begin{columns}
\column{.5\textwidth}
\begin{itemize}
\item Testamos se há elementos;
\item Decrementamos o fim da fila (último);
\item Salvamos o primeiro elemento em variável auxiliar;
\item Empurramos tudo para a frente.
\end{itemize}
\column{.5\textwidth}
\begin{center}
\begin{tabular}{| p{.5cm} |p{.5cm}c }
\cline{1-1}
 \cellcolor{SpringGreen} 24 & &\\ \cline{1-1}
\cellcolor{GreenYellow}  89& &\\ \cline{1-1}
\cellcolor{SpringGreen}  12& &\\ \cline{1-1}
\cellcolor{GreenYellow}  4 & &\\ \cline{1-1}
\cellcolor{SpringGreen} 55 & &\\ \cline{1-2}
\cellcolor{GreenYellow}  & \multicolumn{1}{|p{.5cm}|} {4} & \\ \cline{1-2}
\end{tabular}
\begin{tabular}{| p{.5cm} | }
\hline
\cellcolor{Mahogany} {20} \\\hline
\end{tabular}
\end{center}
\end{columns}
\end{frame}

\begin{frame}[fragile]{Método \textit{T retira()}}
\begin{columns}
\column{.5\textwidth}
\begin{itemize}
\item Testamos se há elementos;
\item Decrementamos o fim da fila (último);
\item Salvamos o primeiro elemento em variável auxiliar;
\item Empurramos tudo para a frente.
\end{itemize}
\column{.5\textwidth}
\begin{center}
\begin{tabular}{| p{.5cm} |p{.5cm}c }
\cline{1-1}
 \cellcolor{SpringGreen} 24 & &\\ \cline{1-1}
\cellcolor{GreenYellow}  89& &\\ \cline{1-1}
\cellcolor{SpringGreen}  12& &\\ \cline{1-1}
\cellcolor{GreenYellow}  4 & &\\ \cline{1-1}
\cellcolor{SpringGreen}  & &\\ \cline{1-2}
\cellcolor{GreenYellow} 55 & \multicolumn{1}{|p{.5cm}|} {4} & \\ \cline{1-2}
\end{tabular}
\begin{tabular}{| p{.5cm} | }
\hline
\cellcolor{Mahogany} {20} \\\hline
\end{tabular}
\end{center}
\end{columns}
\end{frame}

\begin{frame}[fragile]{Método \textit{T retira()}}
\begin{columns}
\column{.5\textwidth}
\begin{itemize}
\item Testamos se há elementos;
\item Decrementamos o fim da fila (último);
\item Salvamos o primeiro elemento em variável auxiliar;
\item Empurramos tudo para a frente.
\end{itemize}
\column{.5\textwidth}
\begin{center}
\begin{tabular}{| p{.5cm} |p{.5cm}c }
\cline{1-1}
 \cellcolor{SpringGreen} 24 & &\\ \cline{1-1}
\cellcolor{GreenYellow}  89& &\\ \cline{1-1}
\cellcolor{SpringGreen}  12& &\\ \cline{1-1}
\cellcolor{GreenYellow}   & &\\ \cline{1-1}
\cellcolor{SpringGreen} 4 & &\\ \cline{1-2}
\cellcolor{GreenYellow} 55 & \multicolumn{1}{|p{.5cm}|} {4} & \\ \cline{1-2}
\end{tabular}
\begin{tabular}{| p{.5cm} | }
\hline
\cellcolor{Mahogany} {20} \\\hline
\end{tabular}
\end{center}
\end{columns}
\end{frame}

\begin{frame}[fragile]{Método \textit{T retira()}}
\begin{columns}
\column{.5\textwidth}
\begin{itemize}
\item Testamos se há elementos;
\item Decrementamos o fim da fila (último);
\item Salvamos o primeiro elemento em variável auxiliar;
\item Empurramos tudo para a frente.
\end{itemize}
\column{.5\textwidth}
\begin{center}
\begin{tabular}{| p{.5cm} |p{.5cm}c }
\cline{1-1}
 \cellcolor{SpringGreen} 24 & &\\ \cline{1-1}
\cellcolor{GreenYellow}  89& &\\ \cline{1-1}
\cellcolor{SpringGreen}  & &\\ \cline{1-1}
\cellcolor{GreenYellow}  12 & &\\ \cline{1-1}
\cellcolor{SpringGreen} 4 & &\\ \cline{1-2}
\cellcolor{GreenYellow} 55 & \multicolumn{1}{|p{.5cm}|} {4} & \\ \cline{1-2}
\end{tabular}
\begin{tabular}{| p{.5cm} | }
\hline
\cellcolor{Mahogany} {20} \\\hline
\end{tabular}
\end{center}
\end{columns}
\end{frame}

\begin{frame}[fragile]{Método \textit{T retira()}}
\begin{columns}
\column{.5\textwidth}
\begin{itemize}
\item Testamos se há elementos;
\item Decrementamos o fim da fila (último);
\item Salvamos o primeiro elemento em variável auxiliar;
\item Empurramos tudo para a frente.
\end{itemize}
\column{.5\textwidth}
\begin{center}
\begin{tabular}{| p{.5cm} |p{.5cm}c }
\cline{1-1}
 \cellcolor{SpringGreen} 24 & &\\ \cline{1-1}
\cellcolor{GreenYellow}  & &\\ \cline{1-1}
\cellcolor{SpringGreen}  89& &\\ \cline{1-1}
\cellcolor{GreenYellow}  12 & &\\ \cline{1-1}
\cellcolor{SpringGreen} 4 & &\\ \cline{1-2}
\cellcolor{GreenYellow} 55 & \multicolumn{1}{|p{.5cm}|} {4} & \\ \cline{1-2}
\end{tabular}
\begin{tabular}{| p{.5cm} | }
\hline
\cellcolor{Mahogany} {20} \\\hline
\end{tabular}
\end{center}
\end{columns}
\end{frame}

\begin{frame}[fragile]{Método \textit{T retira()}}
\begin{columns}
\column{.5\textwidth}
\begin{itemize}
\item Testamos se há elementos;
\item Decrementamos o fim da fila (último);
\item Salvamos o primeiro elemento em variável auxiliar;
\item Empurramos tudo para a frente.
\end{itemize}
\column{.5\textwidth}
\begin{center}
\begin{tabular}{| p{.5cm} |p{.5cm}c }
\cline{1-1}
 \cellcolor{SpringGreen}  & &\\ \cline{1-1}
\cellcolor{GreenYellow}  24& &\\ \cline{1-1}
\cellcolor{SpringGreen}  89& &\\ \cline{1-1}
\cellcolor{GreenYellow}  12 & &\\ \cline{1-1}
\cellcolor{SpringGreen} 4 & &\\ \cline{1-2}
\cellcolor{GreenYellow} 55 & \multicolumn{1}{|p{.5cm}|} {4} & \\ \cline{1-2}
\end{tabular}
\begin{tabular}{| p{.5cm} | }
\hline
\cellcolor{Mahogany} {20} \\\hline
\end{tabular}
\end{center}
\end{columns}
\end{frame}

\begin{frame}[fragile]{Trabalho Fila de Vetor}
\lstset{language=C++,
          keywordstyle=\color{blue}\ttfamily,
          stringstyle=\color{red}\ttfamily,
          commentstyle=\color{OliveGreen}\ttfamily,
          breaklines=true,
          basicstyle=\ttfamily\footnotesize
          }
\begin{itemize}
\item Implemente uma classe Fila todas as operações vistas;
\item Implemente a fila usando Templates;
\item Implemente a fila com um numero de elementos variável definido na instanciação;
\item Use as melhores práticas de orientação a objetos;
\item Documente todas as classes, métodos e atributos;
\item Aplique os testes unitários disponíveis no moodle da disciplina para validar sua estrutura de dados;
\item Entregue até a data definida no moodle.
\end{itemize}
\end{frame}

{
\usebackgroundtemplate{\includegraphics[width=\paperwidth,
height=\paperheight]{../reusable_images/fundo_capa.png}}
\begin{frame}

{\LARGE Perguntas????}

\end{frame}
}


{
\usebackgroundtemplate{\includegraphics[width=\paperwidth,
height=\paperheight]{../reusable_images/fundo_capa.png}}
\begin{frame}
\includegraphics[scale=0.8]{../reusable_images/cc_logo_arge.png}\hspace{0.5cm}
\includegraphics[scale=0.95]{../reusable_images/by.png}

\vspace{1cm}
Este trabalho está licenciado sob uma Licença Creative Commons Atribuição 4.0 Internacional. Para ver uma cópia desta licença, visite http://creativecommons.org/licenses/by/4.0/.

\end{frame}
}
\end{document}
